\documentclass{beamer}
%
% Choose how your presentation looks.
%
% For more themes, color themes and font themes, see:
% http://deic.uab.es/~iblanes/beamer_gallery/index_by_theme.html
%
\mode<presentation>
{
  \usetheme{Singapore}      % or try Darmstadt, Madrid, Warsaw, ...
  \usecolortheme{crane} % or try albatross, beaver, crane, ...
  \usefonttheme{serif}  % or try serif, structurebold, ...
  \setbeamertemplate{navigation symbols}{}
  \setbeamertemplate{caption}[numbered]
} 

\usepackage[english]{babel}
\usepackage[autostyle=true]{csquotes}
\usepackage{xeCJK}
\usepackage[T1]{fontenc}
\usepackage{xcolor}
\setmainfont{Charis SIL}
\setCJKmainfont{Noto Serif CJK HK}

\title[Pearl River Sinitic]{Sinitic Languages of the Pearl River Region \\ --- Unit 1}
\subtitle{13\textsuperscript{th} EACL Summer School, Universität Stuttgart}
\author{Hilário de Sousa 蘇沙}
\institute{EHESS–CRLAO \\ membre associé}
\date{Monday 27 July 2026 \\ 丙午六月十四 \\ v 1.0}

\begin{document}

\begin{frame}
  \titlepage
\end{frame}

% Uncomment these lines for an automatically generated outline.
%\begin{frame}{Outline}
%  \tableofcontents
%\end{frame}

\section{Preliminary matters}

\begin{frame}{Preliminary matters}

\begin{itemize}
  \item Lecturer: Dr. Hilário de Sousa \textsuperscript{\emph{any animate pronoun}} \\ (First-name basis OK: /hɪˈlaːɹɪɐʉ/ /hiˈlɛːɹɪɐʉ/ /iˈlaɾiu/ /ilaʁjo/ /ɦiˈlaːɾijoː/ /hiˈlaːʁio/ etc., \\ or \emph{蘇沙 sou1 saa4} in Cantonese)
  \item Written communications in Chinese can be made in 繁體字 or 简体字
  \item Sometimes displayed as 繁[简], e.g. \\ 鬆[松] \emph{sung1} \enquote*{loose}, 松[松] \emph{cung4} \enquote*{pine} \\ 裏/裡[里] \emph{leoi5} \enquote*{inside}, 里[里] \emph{lei5} \enquote*{neighbourhood}
\end{itemize}

\end{frame}

\begin{frame}{Preliminary matters}

\begin{itemize}
  \item Short clarification questions can be asked at any time; for longer questions, wait till I ask \enquote{are there any questions?} or when I announce \enquote{next slide}
  \item Sharing of academic knowledge is welcome (but do not over-share)
  \item I aim to have a 5-ish-minute break in the middle of each lecture
  
\end{itemize}

\end{frame}

\begin{frame}{Preliminary matters}

\begin{itemize}
  
  \item There is no course-work in this course; contact me if your institution requires you to do course-work for this course
  \item Slides can be downloaded from my website \url{https://hilario.bambooradical.com} \\ \quad (bamboo radical: ⺮\ ) \\ \quad \emph{check for the latest version before each lecture}
\end{itemize}

\end{frame}

\section{Introduction}

\begin{frame}[fragile]{Introduction --- Pearl River Basin}

\begin{figure}
\includegraphics[width=0.8\textwidth]{pearlriver}
\caption{\label{fig:pearlriver}Pearl River Basin}
\end{figure}

\vfill {\tiny Based on \url{https://en.wikipedia.org/wiki/Pearl_River#/media/File:Zhujiangrivermap.png}}

%\begin{block}{Examples}
%Some examples of commonly used commands and features are included, %to help you get started.
%\end{block}

\end{frame}

\begin{frame}[fragile]{Introduction --- Sinitic languages}

漢語方言 \emph{hon3 jyu5 fong1 jin4 / hànyǔ fāngyán}: \\ \enquote{Sinitic languages} / \enquote{Chinese dialects} 

\\\

「方言」 partially overlaps with \enquote{dialect} \\
方言：regional speech variety (language or dialect) \\
\quad \emph{usually used in opposition to a standard speech variety} \\
\textcolor{black!20}{dialect: speech varieties of a location\\
\quad language: collection of \enquote{mutually intelligible} dialects}

\end{frame}

\begin{frame}[fragile]{Introduction --- Sinitic languages}

\begin{figure}
\centering
\includegraphics[width=1.1\textwidth]{mandarinisalanguage}
\caption{\label{fig:mandarinisalanguage}language/dialect as direct translation of 语言/方言}
\end{figure}

\vfill {\tiny Malaysian English from YouTube video \verb+NOODOU教你讲新宁话! Guess The Phrases in Toishanese/Taishanese/Sinneng+ \url{https://youtu.be/cmR53pHTLJI}}

\end{frame}

\begin{frame}[fragile]{Introduction --- Sinitic languages}

漢語方言 \emph{hon3 jyu5 fong1 jin4 / hànyǔ fāngyán}: \\ \enquote{Sinitic languages} / \enquote{Chinese dialects} 

\\\

「方言」 overlaps with \enquote{dialect} \\
\textcolor{black!20}{方言：regional speech variety (language or dialect) \\
\quad \emph{usually used in opposition to a standard speech variety}} \\
dialect: speech variety of a location\\
\quad language: collection of \enquote{mutually intelligible} dialects

\end{frame}

\begin{frame}[fragile]{Introduction --- Sinitic languages}

\begin{figure}
\centering
\includegraphics[width=0.8\textwidth]{germandialects}
\caption{\label{fig:germandialects} \enquote{German dialects}}
\end{figure}

\end{frame}


\begin{frame}[fragile]{Introduction --- Sinitic languages}

Sinitic languages \textsuperscript{purely linguistic term}: \\ Old Chinese and its descendants \\ 
\emph{Old Chinese, e.g.: \\
\quad Oracle Bone Script 甲骨文 (13\textsuperscript{th}–11\textsuperscript{th} C BCE) \\ \quad 4 Books 5 Classics 四書五經 (before 2\textsuperscript{nd} C BCE)}

\\\

\enquote{Chinese}:
\begin{itemize}
    \item Country of China
    \item Chinese dynasties / regimes
    \item Hàn 漢 (Tong4 唐) ethnicity
    \item etc.
\end{itemize}

\end{frame}

\begin{frame}[fragile]{Introduction --- Sinitic languages}

\begin{enumerate}
    \item \enquote{Chinese dialects}/\enquote{Chinese languages} can be interpreted as \enquote{dialects/languages of China} \emph{(which can be Sinitic or non-Sinitic)}
    \item Communities of native speakers of Sinitic languages can identify with all, some, or none of the non-linguistic senses of \enquote{Chinese}, 
    \emph{e.g. Sán Dìu 山瑤 of Vietnam identifies with none}

\centering\includegraphics[width=0.50\textwidth]{clbsandiu1}

\end{enumerate}

\end{frame}

\begin{frame}[fragile]{Introduction --- Sinitic languages}

\begin{figure}
\includegraphics[width=0.55\textwidth]{lacdhuangetal2024}
\caption{\label{fig:LACDRGB}RGB–Multidimensional Scaling map of Southeastern Sinitic}
\end{figure}

\vfill {\tiny Huang, He \& Grieve, Jack \& Jiao, Lei \& Cai, Zhuo. 2024. Geographic structure of Chinese dialects: a computational dialectometric approach. \emph{Linguistics} 62(4). 937–976. \url{https://doi.org/10.1515/ling-2021-0138} \\ Based on data from the \emph{Linguistic Atlas of Chinese Dialects 汉语方言地图集}}

\end{frame}

\begin{frame}[fragile]{Introduction --- Sinitic languages}

\begin{figure}
\includegraphics[width=0.55\textwidth]{lacdhuangetal2024cont}
\caption{\label{fig:LACDRGBcont}Some Sinitic dialect continua}
\end{figure}

\vfill {\tiny Base map from Huang et al. (2024), based on 510 features in the \emph{LACD}}

\end{frame}

\begin{frame}[fragile]{Introduction --- Sinitic languages}

\emph{Language Atlas of China}: 10 Sinitic dialect groups \textsuperscript{(much debated!)} \\ {\scriptsize 官話 Mandarin, 晉 Jìn, 湘 Xiāng, 贛 Gàn, 徽 Huī, 吳 Wú, 閩 Mǐn, 客家 Hakka, 粵 Yuè, 平話 Pínghuà (+ smaller unclassified groups)} 

\begin{figure}
\includegraphics[width=0.4\textwidth]{southsinitic}
\caption{\label{fig:pearlriver}Southern Sinitic dialect groups}
\end{figure}

\vfill \tiny \url{https://en.wikipedia.org/wiki/Sinitic_languages#/media/File:Map_of_sinitic_languages_full-en.svg} Based on 1\textsuperscript{st} ed. of the \emph{Language Atlas of China 中国语言地图集}

\end{frame}


\section{乜}

\subsection{Tables and Figures}

\begin{frame}{Tables and Figures}

\begin{itemize}
\item Use \texttt{tabular} for basic tables --- see Table~\ref{tab:widgets}, for example.
\item You can upload a figure (JPEG, PNG or PDF) using the files menu. 
\item To include it in your document, use the \texttt{includegraphics} command (see the comment below in the source code).
\end{itemize}

\begin{table}
\centering
\begin{tabular}{l|r}
Item & Quantity \\\hline
Widgets & 42 \\
Gadgets & 13
\end{tabular}
\caption{\label{tab:widgets}An example table.}
\end{table}

\end{frame}

\subsection{Mathematics}

\begin{frame}{Readable Mathematics}

Let $X_1, X_2, \ldots, X_n$ be a sequence of independent and identically distributed random variables with $\text{E}[X_i] = \mu$ and $\text{Var}[X_i] = \sigma^2 < \infty$, and let
\[ S_n = \frac{X_1 + X_2 + \cdots + X_n}{n}
      = \frac{1}{n}\sum_{i}^{n} X_i \]
denote their mean. Then as $n$ approaches infinity, the random variables $\sqrt{n}(S_n - \mu)$ converge in distribution to a normal $\mathcal{N}(0, \sigma^2)$.

\end{frame}

\end{document}
